% SPDX-License-Identifier: CC-BY-SA-4.0
% Author: Matthieu Perrin

\begingroup

\begin{exercice}[Révisions -- Preuves mathématiques]\label{exo:introduction/maths/revision}

  Démontrer les propriétés suivantes sur les ensembles. 
  \begin{question}
  \item $\forall A, \forall B, A \subseteq A \cup B$.
  \item $\forall A, \forall B, (A \setminus B) \cap B = \emptyset$.
  \item $\forall A, \forall B, \forall C, A \cup (B \cap C) = (A \cup B) \cap (A \cup C)$.
  \item Montrer par récurrence que pour tout entier naturel $n$, 
    $$
    \sum_{i=1}^n i = 1 + 2 + \dots + n = \frac{n(n+1)}{2}.
    $$
  \item Montrer par récurrence que pour tout entier naturel $n$,
    $$
    \sum_{i=0}^n 2^i = 1 + 2 + 4 + \dots + 2^n = 2^{n+1} - 1.
    $$
  \item On définit la suite $(u_n)_{n\in\mathbb{N}}$ par $u_0 = 1$ et $u_{n+1} = 3u_n + 2$.
    Montrer par récurrence que pour tout $n$,
    $$
    u_n = 2 \cdot 3^n - 1.
    $$
  \end{question}
  
\end{exercice}

\begin{correction}

  \begin{question}

  \item Montrer que $\forall A, \forall B, \; A \subseteq A \cup B$.
    \begin{itemize}
    \item Soit $x \in A$. Par définition de $\cup$, $x \in A \cup B$. Donc $A \subseteq A \cup B$.
    \end{itemize}

  \item Montrer que $\forall A, \forall B, \; (A \setminus B) \cap B = \emptyset$.
    \begin{itemize}
    \item Raisonnons par l’absurde. Supposons qu’il existe $x \in (A \setminus B) \cap B$.
    \item Alors $x \in A \setminus B$ et $x \in B$. Par définition de $\setminus$, $x \in A$ et $x \notin B$, contradiction avec $x \in B$.
    \item Donc $(A \setminus B) \cap B = \emptyset$.
    \end{itemize}

  \item Montrer que $\forall A, \forall B, \forall C, \; A \cup (B \cap C) = (A \cup B) \cap (A \cup C)$.
    \begin{itemize}
    \item \textit{Sens $\subseteq$ :} Soit $x \in A \cup (B \cap C)$. Alors $x \in A$ ou $x \in B \cap C$.
      \begin{itemize}
      \item Si $x \in A$, alors $x \in A \cup B$ et $x \in A \cup C$, donc $x \in (A \cup B) \cap (A \cup C)$.
      \item Si $x \in B \cap C$, alors $x \in B$ et $x \in C$, donc $x \in A \cup B$ et $x \in A \cup C$, donc $x \in (A \cup B) \cap (A \cup C)$.
      \end{itemize}
    \item \textit{Sens $\supseteq$ :} Soit $x \in (A \cup B) \cap (A \cup C)$. Alors $x \in A \cup B$ et $x \in A \cup C$.
      \begin{itemize}
      \item Si $x \in A$, alors $x \in A \cup (B \cap C)$.
      \item Sinon $x \notin A$. Comme $x \in A \cup B$, on a $x \in B$. Comme $x \in A \cup C$, on a $x \in C$. Donc $x \in B \cap C$, donc $x \in A \cup (B \cap C)$.
      \end{itemize}
    \end{itemize}
    Par double inclusion, $A \cup (B \cap C) = (A \cup B) \cap (A \cup C)$.

  \item Preuve que $\forall n \in \mathbb{N}, \; \sum_{i=1}^n i = \frac{n(n+1)}{2}$.
    \begin{description}
    \item[Initialisation :] Pour $n=0$ : $0 = \frac{0 \times 1}{2}$.
    \item[Hérédité :] Supposons la propriété vraie au rang $n$ : $\sum_{i=1}^n i = \frac{n(n+1)}{2}$.
      Alors
      $$
      \sum_{i=1}^{n+1} i = \left(\sum_{i=1}^n i\right) + (n+1)
      = \frac{n(n+1)}{2} + (n+1)
      = (n+1)\left(\frac{n}{2}+1\right)
      = \frac{(n+1)(n+2)}{2}.
      $$
    \item[Conclusion :] La propriété est vraie pour tout $n \in \mathbb{N}$.
    \end{description}

  \item Preuve que $\forall n \in \mathbb{N}, \; \sum_{i=0}^n 2^i = 2^{n+1} - 1$.
    \begin{description}
    \item[Initialisation :] Pour $n=0$ : $2^0 = 1 = 2^{1} - 1$.
    \item[Hérédité :] Supposons la propriété vraie au rang $n$ : $\sum_{i=0}^n 2^i = 2^{n+1} - 1$.
      Alors
      $$
      \sum_{i=0}^{n+1} 2^i = \left(\sum_{i=0}^{n} 2^i\right) + 2^{n+1}
      = (2^{n+1}-1) + 2^{n+1}
      = 2^{n+2} - 1.
      $$
    \item[Conclusion :] La propriété est vraie pour tout $n \in \mathbb{N}$.
    \end{description}

  \item Soit la suite définie par $u_0 = 1$ et $u_{n+1} = 3u_n + 2$. Montrer que $\forall n, \; u_n = 2 \cdot 3^n - 1$.
    \begin{description}
    \item[Initialisation :] Pour $n=0$ : $u_0 = 1$ et $2 \cdot 3^0 - 1 = 1$.
    \item[Hérédité :] Supposons que $u_n = 2 \cdot 3^n - 1$ pour un certain $n$. Alors
      $$
      u_{n+1} = 3u_n + 2
      = 3(2 \cdot 3^n - 1) + 2
      = 2 \cdot 3^{n+1} - 3 + 2
      = 2 \cdot 3^{n+1} - 1.
      $$
    \item[Conclusion :] La formule est vraie pour tout $n \in \mathbb{N}$.
    \end{description}

  \end{question}

\end{correction}

\endgroup
\endinput
